\input{../tex-inputs/latex-korrekturansicht-vorspann}

\section[Arthur Schnitzler an Berta Zuckerkandl, 9. 12. 1911]{L03951 Arthur Schnitzler an Berta Zuckerkandl, 9. 12. 1911}
\nopagebreak\mylabel{L03951v}
\rehead{ }\normalsize\beginnumbering\briefempfaengerindex{Zuckerkandl, Berta@\textsc{Zuckerkandl, Berta}!zzzSchnitzler, Arthur@\emph{von Arthur Schnitzler}!1911-12-091@{9. 12. 1911}|(be}
\toendnotes[C]{\smallbreak\pagebreak[2]}
\correspDesc{Versand  durch Arthur Schnitzler am 9. 12. 1911 in Wien
\newline{}Erhalt  durch Berta Zuckerkandl im Zeitraum [9. 12. 1911
                  – 12. 12. 1911?] in Wien}\toendnotes[C]{\smallbreak}
\Standort{DLA, HS.1985.1.2282.}
\physDesc{Brief, Durchschlag, 1 Blatt, 2 Seiten, 1069 Zeichen
\newline{}Schreibmaschine
\newline{}Handschrift: roter Buntstift, lateinische Kurrent (\noindent{}beschriftet: »\uline{\textsc{Zuckerkandl}}« und »\textsc{Fr}{[}an{]}\textsc{k}{[}reich{]}«, vier Unterstreichungen)}\toendnotes[C]{\smallbreak}
\pstart
           \raggedleft{}{\pb}9. 12. 1911.\pend
           
\pstart\center{}Verehrte gnädige Frau.\pend\vspace{0.5em}
\pstart
           Herzlichen Dank für die freundliche Uebersendung des \label{K_L03951-1v}\edtext{Briefes von Madame \textcolor{blue}{Lefèvre}\pwindex{Lefèvre, A. @\textsc{Lefèvre, A.}, \emph{Übersetzerin}|pw}{}\ledrightnote{\textcolor{blue}{A. Lefèvre}}}{\lemma{\textnormal{\emph{Briefes … Lefèvre}}}\Cendnote{\textnormal{nicht
                  überliefert}}}\label{K_L03951-1}.
               Ich habe also gleich \label{K_L03951-2v}\edtext{an sie
                  geschrieben}{\lemma{\textnormal{\emph{an sie
                  geschrieben}}}\Cendnote{\textnormal{Es sind keine Briefe \textcolor{blue}{Schnitzlers} an \textcolor{blue}{Levèvre}\pwindex{Lefèvre, A. @\textsc{Lefèvre, A.}, \emph{Übersetzerin}|pwk} überliefert.}}}\label{K_L03951-2} und sie vor allem um
               Beantwortung der Frage ersucht, ob\textcolor{blue}{Guitry}\pwindex{Guitry, Lucien 13.\,12.\,1860 Paris – 1.\,6.\,1925 ebd.@\textsc{Guitry, Lucien} (13.\,12.\,1860 Paris – 1.\,6.\,1925 ebd.), \emph{Schriftsteller, Schauspieler}|pw}{}\ledrightnote{\textcolor{blue}{Lucien Guitry}} im nächsten Jahre bei \textcolor{blue}{Antoine}\pwindex{Antoine, André 31.\,1.\,1858 Limoges – 23.\,10.\,1943 Le Pouliguen@\textsc{Antoine, André} (31.\,1.\,1858 Limoges – 23.\,10.\,1943 Le Pouliguen), \emph{Theaterleiter, Schauspieler}|pw}{}\ledrightnote{\textcolor{blue}{André Antoine}} sein wird, ferner ob sie überzeugt ist, dass \textcolor{blue}{Antoine}\pwindex{Antoine, André 31.\,1.\,1858 Limoges – 23.\,10.\,1943 Le Pouliguen@\textsc{Antoine, André} (31.\,1.\,1858 Limoges – 23.\,10.\,1943 Le Pouliguen), \emph{Theaterleiter, Schauspieler}|pw}{}\ledrightnote{\textcolor{blue}{André Antoine}} ihr eine Uebersetzung besonders gerne
               anvertrauen würde. Denn schliesslich Herr \textcolor{blue}{Remon}\pwindex{Rémon, Maurice 27.\,11.\,1861 Paris – 20.\,6.\,1945 Mérignac@\textsc{Rémon, Maurice} (27.\,11.\,1861 Paris – 20.\,6.\,1945 Mérignac), \emph{Übersetzer}|pw}{}\ledrightnote{\textcolor{blue}{Maurice Rémon}} hat mir über seine Verbindung mit \textcolor{blue}{Antoine}\pwindex{Antoine, André 31.\,1.\,1858 Limoges – 23.\,10.\,1943 Le Pouliguen@\textsc{Antoine, André} (31.\,1.\,1858 Limoges – 23.\,10.\,1943 Le Pouliguen), \emph{Theaterleiter, Schauspieler}|pw}{}\ledrightnote{\textcolor{blue}{André Antoine}} ungefähr \label{K_L03951-3v}\edtext{dasselbe
                  geschrieben}{\lemma{\textnormal{\emph{dasselbe
                  geschrieben}}}\Cendnote{\textnormal{\textcolor{blue}{Maurice Rémon}\pwindex{Rémon, Maurice 27.\,11.\,1861 Paris – 20.\,6.\,1945 Mérignac@\textsc{Rémon, Maurice} (27.\,11.\,1861 Paris – 20.\,6.\,1945 Mérignac), \emph{Übersetzer}|pwk} an \textcolor{blue}{Arthur Schnitzler}, 1. 11. 1911, vgl.
                     Karl Zieger: \emph{Arthur Schnitzler et la France 1894–1938.
                        Enquête sur une réception}, Villeneuve d’Ascq:
                        \emph{Presses Universitaires du
                        Septentrion} 2012, S. 190.
                  }}}\label{K_L03951-3} und behauptet noch überdies,
               dass seine \textcolor{green}{Uebersetzung}\pwindex{Schnitzler, Arthur 15. 5. 1862 Wien – 21. 10. 1931 ebd.@\textsc{Schnitzler, Arthur} (15. 5. 1862 Wien – 21. 10. 1931 ebd.), \emph{Schriftsteller, Mediziner}!femme au poignard@\strich\emph{La femme au poignard}|pwv}{}\ledrightnote{{$\rightarrow$}\emph{\textcolor{green}{La femme au poignard}}} der
                  »\textcolor{green}{Frau mit dem Dolch}\pwindex{Schnitzler, Arthur 15. 5. 1862 Wien – 21. 10. 1931 ebd.@\textsc{Schnitzler, Arthur} (15. 5. 1862 Wien – 21. 10. 1931 ebd.), \emph{Schriftsteller, Mediziner}!Frau mit dem Dolche@\strich\emph{Die Frau mit dem Dolche}|pw}{}\ledrightnote{\textcolor{green}{Die Frau mit dem Dolche}}« heuer bei \textcolor{blue}{Antoine}\pwindex{Antoine, André 31.\,1.\,1858 Limoges – 23.\,10.\,1943 Le Pouliguen@\textsc{Antoine, André} (31.\,1.\,1858 Limoges – 23.\,10.\,1943 Le Pouliguen), \emph{Theaterleiter, Schauspieler}|pw}{}\ledrightnote{\textcolor{blue}{André Antoine}} zur Aufführung kommen soll. Nebstbei
               scheint er doch in direkter Verbindung mit \textcolor{blue}{Guitry}\pwindex{Guitry, Lucien 13.\,12.\,1860 Paris – 1.\,6.\,1925 ebd.@\textsc{Guitry, Lucien} (13.\,12.\,1860 Paris – 1.\,6.\,1925 ebd.), \emph{Schriftsteller, Schauspieler}|pw}{}\ledrightnote{\textcolor{blue}{Lucien Guitry}} zu stehen und es kommt in der Angelegenheit des »\textcolor{green}{Weiten Landes}\pwindex{Schnitzler, Arthur 15. 5. 1862 Wien – 21. 10. 1931 ebd.@\textsc{Schnitzler, Arthur} (15. 5. 1862 Wien – 21. 10. 1931 ebd.), \emph{Schriftsteller, Mediziner}!weite Land. Tragikomödie in fünf Akten@\strich\emph{Das weite Land. Tragikomödie in fünf Akten}|pw}{}\ledrightnote{\textcolor{green}{Das weite Land. Tragikomödie in fünf Akten}}« doch mehr auf \textcolor{blue}{Guitry}\pwindex{Guitry, Lucien 13.\,12.\,1860 Paris – 1.\,6.\,1925 ebd.@\textsc{Guitry, Lucien} (13.\,12.\,1860 Paris – 1.\,6.\,1925 ebd.), \emph{Schriftsteller, Schauspieler}|pw}{}\ledrightnote{\textcolor{blue}{Lucien Guitry}} an als auf \textcolor{blue}{Antoine}\pwindex{Antoine, André 31.\,1.\,1858 Limoges – 23.\,10.\,1943 Le Pouliguen@\textsc{Antoine, André} (31.\,1.\,1858 Limoges – 23.\,10.\,1943 Le Pouliguen), \emph{Theaterleiter, Schauspieler}|pw}{}\ledrightnote{\textcolor{blue}{André Antoine}}, umso mehr als \textcolor{blue}{Antoine}\pwindex{Antoine, André 31.\,1.\,1858 Limoges – 23.\,10.\,1943 Le Pouliguen@\textsc{Antoine, André} (31.\,1.\,1858 Limoges – 23.\,10.\,1943 Le Pouliguen), \emph{Theaterleiter, Schauspieler}|pw}{}\ledrightnote{\textcolor{blue}{André Antoine}} es
               nicht einmal der Mühe wert hält mir auf meinen direkten und klare Fragen
               formulierenden \label{K_L03951-4v}\edtext{Brief}{\lemma{\textnormal{\emph{Brief}}}\Cendnote{\textnormal{\textcolor{blue}{Arthur Schnitzler} an \textcolor{blue}{André Antoine}\pwindex{Antoine, André 31.\,1.\,1858 Limoges – 23.\,10.\,1943 Le Pouliguen@\textsc{Antoine, André} (31.\,1.\,1858 Limoges – 23.\,10.\,1943 Le Pouliguen), \emph{Theaterleiter, Schauspieler}|pwk}, 20. 11. 1911, \emph{Deutsches Literaturarchiv Marbach},
                  HS.1985.1.257.}}}\label{K_L03951-4} zu antworten. {\pb}Wir wollen
               nun die Antwort der Frau \textcolor{blue}{Lefèvre}\pwindex{Lefèvre, A. @\textsc{Lefèvre, A.}, \emph{Übersetzerin}|pw}{}\ledrightnote{\textcolor{blue}{A. Lefèvre}} abwarten und
               dann weiter sehen.\pend
           
\pstart
           Mit wiederholtem Dank für Ihre Bemühungen und Ihr liebenswürdiges Interesse mit
               vielen{\\[\baselineskip]} Grüssen{\\[\baselineskip]} Ihr aufrichtig ergebener\pend
           \leftskip=0em{}{\vspace{1\baselineskip}}
\pstart
           \noindent{}Frau Berta Zuckerkandl, \textcolor{pink}{Wien}\oindex{Wien@\textbf{Wien}, \emph{Verwaltungsgebiet}|pw}{}\ledrightnote{\textcolor{pink}{Wien}}.\pend
           \selectlanguage{ngerman}\endnumbering\briefempfaengerindex{Zuckerkandl, Berta@\textsc{Zuckerkandl, Berta}!zzzSchnitzler, Arthur@\emph{von Arthur Schnitzler}!1911-12-091@{9. 12. 1911}|)be}\mylabel{L03951h}  \newcommand{\dateiname}{L03951}\newcommand{\titel}{Arthur Schnitzler an Berta Zuckerkandl, 9. 12. 1911}\newcommand{\editorInnen}{Herausgegeben von Jahnke, SelmaMüller, Martin Anton}\input{../tex-inputs/latex-korrekturansicht-abspann}
